\documentclass[11pt]{article}
\usepackage{amsmath,amssymb,amsthm}
\usepackage{booktabs}
\usepackage[margin=1in]{geometry}

\newtheorem{theorem}{Theorem}
\newtheorem{lemma}[theorem]{Lemma}

\title{Problem 702: Jumping Flea}
\author{Project Euler Solutions}
\date{}

\begin{document}
\maketitle

\section{Problem Statement}

A hexagonal table of side $N$ is divided into unit triangles. A flea jumps to the midpoint between its current position and a chosen corner. Let $J(T)$ be the minimum jumps needed to reach the interior of triangle $T$, and $S(N)=\sum J(T)$ over upper triangles. Given $S(3)=42$ and $S(5)=126$, find $S(123456789)$.

\section{Mathematical Analysis}

Each jump is a contraction mapping. Position after jumps is determined by the sequence of chosen corners in base-6-like encoding.

\section{Complexity Analysis}

Explicit computation is $O(N^2 \log N)$, or $O(N \log N)$ using the fractal structure.

\section{Answer}

\[
\boxed{622305608172525546}
\]

\end{document}
